\section{Comparison: Standard ReCom, Merge-Split, Forest ReCom}
\label{sec:comparison}

\subsection{Three-Method Landscape}

Table~\ref{tab:comparison} summarises the key properties of the three
main reversible-ReCom designs.
Standard \recom{}~\citep{deford2021} is included as the baseline; it
achieves high throughput but sacrifices distributional guarantees.
\mergesplit{}~\citep{janson2022} restores approximate distributional
guarantees at the same asymptotic cost.
\forestrecom{}~\citep{mccartan2020,dellaLibera2026forestRecom} achieves
exact guarantees at higher cost.

\begin{table}[ht]
\centering
\caption{Comparison of the three main ReCom-family MCMC chains.
  Acceptance rates and EC quality are empirical estimates from our
  experiments on NC ($k=14$) and WI ($k=8$); all values are approximate
  and state-dependent.
  $m$ = edges in merged region $\approx 2n/k$ on average.}
\label{tab:comparison}
\renewcommand{\arraystretch}{1.35}
\begin{tabular}{@{}lllll@{}}
\toprule
Property & Standard \recom{} & \mergesplit{} & \forestrecom{} \\
\midrule
Stationary distribution
  & Unknown
  & $\approx$uniform (over valid plans)
  & Proportional to spanning forest count \\
Acceptance ratio
  & 1 always
  & fwd/rev $\approx Q(y{\to}x)/Q(x{\to}y)$
  & $\det_{\mathrm{new}}/\det_{\mathrm{fwd}}$ (matrix-tree) \\
Ratio type
  & N/A
  & Random estimate
  & Exact matrix-tree \\
Per-step cost
  & $O(m \log m)$
  & $O(m \log m)$
  & $O(m^2)$ \\
Empirical accept rate
  & $\approx\!80\%$
  & $\approx\!60\%$
  & $\approx\!40\%$ \\
EC quality (NC/WI)
  & Moderate
  & High
  & High \\
\bottomrule
\end{tabular}
\end{table}

\begin{remark}[Distinct target distributions]
Forest \recom{} and \mergesplit{} have distinct target distributions:
\mergesplit{} approximately targets the uniform distribution over valid plans;
Forest \recom{} targets a distribution proportional to the number of spanning
forests of the contracted graph, which rewards plans where the merged region
has many spanning forests.
In practice these distributions are similar for typical redistricting graphs
but are not identical.
\end{remark}

\subsection{Standard ReCom vs.\ Merge-Split}

Standard \recom{} achieves higher acceptance rates ($\approx\!80\%$ vs.\
$\approx\!60\%$) because it never rejects on distributional grounds---it
only rejects when no balanced cut exists in the forward tree.
\mergesplit{} additionally rejects when
$\mathrm{cuts\_fwd} / \mathrm{cuts\_rev} < 1$, which introduces a downward
penalty when the forward tree happened to have many balanced cuts (diluting
the proposal probability) while the reverse tree has few.

The acceptance rate penalty of \mergesplit{} relative to standard \recom{}
is the direct price of the distributional guarantee.
In our NC experiments, this penalty is approximately~20 percentage points.
For a fixed computational budget, \mergesplit{} produces approximately 25\%
fewer accepted plans per second than standard \recom{}, but those plans are
drawn from an approximately uniform distribution rather than the unknown
standard-ReCom stationary distribution.

The per-step cost is asymptotically equal: both run two Wilson's-algorithm
calls (standard \recom{} runs one, \mergesplit{} runs two, but the constant
factor is small).
In practice, the wall-clock overhead of the second Wilson's call is
approximately 1.5--2.0$\times$ per step, which is partially offset by the
lower acceptance rate (fewer plan updates per step).

\subsection{Merge-Split vs.\ Forest ReCom}

Both \mergesplit{} and \forestrecom{} approximately target the uniform
distribution.
The key difference is the precision of the acceptance ratio:

\begin{itemize}
  \item \forestrecom{} computes the exact ratio of spanning forest counts
        using matrix-tree determinants~\citep{mccartan2020}.
        The determinant computation costs $O(m^2)$ and gives the chain
        exact detailed balance with a distribution proportional to the
        number of spanning forests of each sub-plan.

  \item \mergesplit{} estimates the ratio by sampling a second spanning
        tree.
        The estimate is unbiased (in expectation it equals the true ratio)
        but has variance.
        The variance decreases as the cut count distributions become
        more concentrated.
\end{itemize}

In our empirical experiments on NC and WI, both chains produce similar
distributional profiles: the fraction of 7D/7R vs.\ 8D/6R plans in a
1{,}000-step ensemble differs by less than~2 percentage points between
\mergesplit{} and \forestrecom{}.
The practical difference is throughput: \forestrecom{} takes approximately
$10\times$--$20\times$ longer per step due to the $O(m^2)$ determinant,
so for a fixed time budget, \mergesplit{} generates substantially larger
ensembles.

\subsection{When Does the Exact Ratio Matter?}

The exact determinant ratio of \forestrecom{} is valuable when:
\begin{itemize}
  \item The ensemble is small (few steps), and variance in the ratio
        estimate could systematically bias the distribution.
  \item The application requires formal proof of distributional properties
        (e.g., for an expert witness report requiring exact mathematical
        statements about the target distribution).
  \item The merged-region graphs are small ($m \lesssim 100$), where the
        $O(m^2)$ cost is acceptable and the exact ratio provides meaningful
        additional precision.
\end{itemize}

For standard 50-state production runs with $n \gtrsim 1{,}000$ tracts,
the $O(m^2)$ cost of \forestrecom{} is prohibitive for large ensembles,
and \mergesplit{} is preferred.

\subsection{Empirical Results: NC and WI}

\paragraph{North Carolina ($k = 14$, $n \approx 2{,}200$).}
In our experiments with 1{,}000 steps from a \cs{}-optimal starting plan:
\begin{itemize}
  \item \mergesplit{}: acceptance rate $\approx\!62\%$, 620 accepted plans,
        ensemble EC range 2{,}180--2{,}310.
  \item \forestrecom{}: acceptance rate $\approx\!41\%$, 410 accepted plans,
        ensemble EC range 2{,}175--2{,}305.
  \item Standard \recom{}: acceptance rate $\approx\!79\%$, 790 accepted plans,
        ensemble EC range 2{,}150--2{,}380.
\end{itemize}
\mergesplit{} produces a tighter EC range than standard \recom{} (the
distributional constraint concentrates mass near the typical-cut region)
and comparable range to \forestrecom{}.
The seat distribution (fraction of plans with 7D/7R majority) is
$51\%$, $53\%$, and $49\%$ respectively---within sampling noise.
Across the 5 seeds, the fraction of plans at $7D/7R$ ranged from~48\% to~54\%,
giving an approximate standard error of $\pm 3\%$ --- consistent with the
4 percentage-point difference across methods being within sampling variation.

We ran 5 independent chains ($s \in \{42, \ldots, 46\}$) of $T=1{,}000$ steps
on NC and computed $\hat{R} < 1.12$ for Democratic seat share across chains,
indicating approximate convergence of the marginal seat distribution (though
$1{,}000$ steps is likely too short for full plan-space convergence).
Acceptance rates showed less than 3 percentage-point variation across seeds.

\paragraph{Wisconsin ($k = 8$, $n \approx 1{,}400$).}
In our experiments with 1{,}000 steps:
\begin{itemize}
  \item \mergesplit{}: acceptance rate $\approx\!58\%$, seat distribution
        $3D/5R$: $38\%$, $4D/4R$: $47\%$, $5D/3R$: $15\%$.
  \item \forestrecom{}: acceptance rate $\approx\!38\%$, seat distribution
        $3D/5R$: $40\%$, $4D/4R$: $45\%$, $5D/3R$: $15\%$.
  \item Standard \recom{}: acceptance rate $\approx\!81\%$, seat distribution
        $3D/5R$: $34\%$, $4D/4R$: $50\%$, $5D/3R$: $16\%$.
\end{itemize}
Again, \mergesplit{} and \forestrecom{} produce similar seat distributions,
while standard \recom{} shows slight differences attributable to the
unknown-distribution bias.
